\documentclass[11pt,a4paper]{article}
\usepackage[margin=1in]{geometry}
\usepackage{hyperref}
\usepackage{enumitem}
\usepackage{graphicx}
\usepackage{xcolor}

% -----------------------------------------------------
% bvraghav-tiet-ucs503p-article.sty
% -----------------------------------------------------

% Variable: Lab Instructor
\makeatletter \newcommand{\BvrSetLabInstructor}[1]{%
  \gdef\Bvr@LabInstructor{#1}}
% default
\newcommand{\Bvr@LabInstructor}{Dr. Raghav %
  B. Venkataramaiyer}
% public accessor
\newcommand{\BvrLabInstructorValue}{\Bvr@LabInstructor}
\makeatother

% Add subtitle
\usepackage{titling}
% -- Set title bold
\pretitle{\begin{center}\bfseries\fontsize{18bp}{18bp}\selectfont}
\makeatletter
\newcommand{\BvrSubtitle}[1]{%
  \posttitle{%
    \par\end{center}
    \begin{center}\large#1\end{center}
    \vskip 0.5em}%
}
\makeatother

% Hints in the section title(s)
\newcommand{\BvrHint}[1]{%
  {\normalfont\normalsize\itshape\hfill #1}}

% -----------------------------------------------------

\title{CreditSetu: A Dual-Path Credit Scoring \& Lending Marketplace for Microfinance}

\BvrSetLabInstructor{Dr. Raghav B. Venkataramaiyer}

\BvrSubtitle{%
  Project Proposal for UCS503P  \\
  Submitted to: \BvrLabInstructorValue \\ \bfseries
  Thapar Institute of Engineering and Technology%
}

\author{
  Maanya Garg, CSED%
  \thanks{Roll No: \texttt{1024030564}, %
    Email: \texttt{<TODO: confirm official Thapar email>}}
  \and
  Kashvi Bansal, CSED%
  \thanks{Roll No: \texttt{1024030563}, %
    Email: \texttt{kbansal3_be24@thapar.edu}}
  \and
  Shreesh Gupta, CSED%
  \thanks{Roll No: \texttt{1024030565}, %
    Email: \texttt{<TODO: confirm official Thapar email>}}
  \and
  Shivam Raj, CSED%
  \thanks{Roll No: \texttt{1024031140}, %
    Email: \texttt{<TODO: confirm official Thapar email>}}
}

\date{\today}

\begin{document}
\maketitle

\section{Higher-order goal%
  \BvrHint{\ldots secondary framing}}
This project aims to widen access to formal credit for
underserved borrowers in India -- particularly members
of rural Self-Help Groups (SHGs) and independent
borrowers who have no prior formal credit history. The
immediate engineering objective is to translate
\emph{trustworthy, explainable credit assessment} into a
measurable, deployable software system that a
microfinance lender could plausibly adopt.

\section{Time-to-value %
\BvrHint{\ldots primary heuristic}}
\begin{itemize}[leftmargin=0.7cm]
    \item We deliver meaningful increments quickly: a
    working synthetic-data pipeline, a trained scoring
    model, and a queryable API were built and exercised
    end-to-end before any UI existed.
    \item We prioritize fast feedback over perfection by
    validating each pipeline stage against real
    (synthetic) data end-to-end -- training on real data,
    scoring real individuals, hitting the live API,
    and screenshotting the live frontend -- rather than
    reading code and assuming correctness.
    \item Success is defined by what can be delivered and
    verified working in each iteration, then improved:
    the scoring engine and open marketplace shipped first;
    role-separated logins and authorization are the
    current iteration.
\end{itemize}

\section{Problem Statement}
Microfinance lending in India serves two kinds of
borrowers very unevenly. Common pain points include:
\begin{itemize}[leftmargin=0.7cm]
    \item \textbf{No credit history for new/independent
    borrowers:} a borrower with no SHG affiliation and no
    prior loans has nothing a lender can score against --
    the classic cold-start problem.
    \item \textbf{Opaque scoring:} where informal scoring
    does exist (e.g. group-level trust), it rarely
    explains \emph{why} a given score is what it is, which
    undermines both borrower trust and lender diligence.
    \item \textbf{No structured SHG-lender relationship:}
    lenders either lend to anyone in a group or nobody --
    there is no formal, auditable approve/reject linkage
    between a specific SHG and a specific lender.
    \item \textbf{Everyone sees everyone's data:} in an
    open system, a borrower, a lender, and an SHG all see
    the same undifferentiated data -- there is no
    role-appropriate, access-controlled view for each
    stakeholder.
\end{itemize}

\textbf{Why this matters now (contextual relevance):} SHG-based
microfinance already operates at large scale in India; a
system that lets a borrower's \emph{own} repayment history
build their score, while still giving SHG members a fair
trust-based head start, directly targets an active
fairness and access gap in that ecosystem.

\section{Proposed Solution %
\BvrHint{\ldots Software Proposition}}
\subsection{Overview}
Build a \textbf{web-based} dual-path credit scoring and
lending-marketplace system, CreditSetu, with the
following components:
\begin{itemize}[leftmargin=0.7cm]
    \item \textbf{Dual-path scoring engine:} an SHG-linked
    borrower's score bootstraps from their group's track
    record; an independent borrower starts from a flat,
    lower base and builds purely from their own repayment
    and savings behaviour.
    \item \textbf{Explainable scoring:} an XGBoost
    classifier with SHAP explanations, so every score
    comes with a feature-level breakdown of why it is what
    it is, not just a number.
    \item \textbf{SHG-lender linkage:} lenders formally
    request, and SHGs approve or reject, a lending
    relationship -- not open lending to anyone.
    \item \textbf{Loan marketplace:} lenders can make
    offers to eligible borrowers, and borrowers can
    directly request loans that are broadcast to every
    lender they are currently eligible for.
    \item \textbf{Role-separated access:} four distinct,
    login-gated views -- Borrower, Lender, SHG, and
    Admin -- each seeing only the data appropriate to that
    role, enforced server-side rather than just hidden in
    the UI.
\end{itemize}

\subsection{Core workflow}
\begin{enumerate}[leftmargin=0.7cm]
    \item A borrower's repayment, savings, and (if
    applicable) SHG-attendance history feeds a feature
    pipeline.
    \item The scoring engine blends a cold-start prior
    (SHG or independent base) with the trained model's
    output as repayment history accumulates.
    \item Lenders discover SHGs and borrowers they are
    formally linked to (or independents, who are open to
    all lenders by design), and either make offers or
    respond to borrower-initiated requests.
    \item Every role logs into its own dashboard: a
    borrower sees only their own score/history/suggestions
    and their own SHG; a lender sees SHGs by area and its
    own linked members and requests; an SHG sees its own
    members and linked lenders; an admin sees read-only
    platform-wide oversight.
\end{enumerate}

\subsection{Operational constraints for an educational
setting}
\begin{itemize}[leftmargin=0.7cm]
    \item Use synthetic but behaviourally realistic data
    (no real borrower PII), generated and versioned as
    part of the pipeline, not hand-curated.
    \item Avoid dependence on novel research algorithms;
    XGBoost + SHAP are mature, well-understood tools --
    the engineering focus is the scoring pipeline, access
    control, and marketplace correctness.
    \item Design for deployability on common web hosting
    (SQLite for the demo, a one-line \texttt{DATABASE\_URL}
    swap to Postgres for concurrent production use).
\end{itemize}

\section{Solution Approach %
\BvrHint{\ldots Engineering focus}}
This is an engineering course project, so we focus on a
correct, explainable, access-controlled pipeline rather
than algorithmic novelty.

\subsection{Scoring pipeline %
\BvrHint{\ldots non-research}}
\begin{itemize}[leftmargin=0.7cm]
    \item Feature engineering turns each individual's
    loans, repayments, savings, and SHG attendance into
    features such as repayment rate, savings regularity,
    and peer repayment rate.
    \item An XGBoost classifier is trained on SHG-linked
    individuals with sufficient repayment history,
    predicting default risk; SHAP's \texttt{TreeExplainer}
    attributes every prediction to its contributing
    features.
    \item A cold-start blend interpolates between a pure
    prior (SHG bootstrap or independent base) and the
    trained model's output as an individual accumulates
    repayment events, so a score is never asked to
    extrapolate from data it doesn't have.
\end{itemize}

\subsection{Web architecture %
\BvrHint{\ldots web-based solution}}
A 3-tier web architecture:
\begin{itemize}[leftmargin=0.7cm]
    \item \textbf{Frontend:} a React (Vite) single-page
    app with four role-gated dashboards and a public
    landing page.
    \item \textbf{Backend API:} FastAPI, with
    bearer-token session auth per role and server-side
    authorization scoping every query to the
    authenticated caller's own identity.
    \item \textbf{Database:} SQLAlchemy over SQLite (demo)
    / Postgres (production), with an append-only score
    history table so a borrower's score-over-time and its
    SHAP explanation are always reconstructable.
\end{itemize}

\subsection{CI/CD and fast delivery enabler}
CI/CD will be used to:
\begin{itemize}[leftmargin=0.7cm]
    \item Automatically build and lint on every change.
    \item Produce a repeatable, from-scratch pipeline run
    (data generation $\to$ training $\to$ batch scoring
    $\to$ serving) as the release artifact, not just a
    code diff.
    \item Enable safe, frequent releases of incremental
    features (the role-separated login redesign is itself
    one such increment, layered on a previously-working
    open system).
\end{itemize}

\section{Evaluation Criterion %
\BvrHint{\ldots measurable and attributable}}
We define success using measurable metrics tied to the
core workflow.

\subsection{Primary evaluation metric}
\textbf{Access-control correctness:} the fraction of
cross-role/cross-identity data-access attempts that are
correctly rejected (401/403) rather than served.
\begin{itemize}[leftmargin=0.7cm]
    \item Target: 100\% of out-of-scope requests (e.g. a
    borrower fetching another borrower's data, or a lender
    fetching an unlinked SHG's member list) are rejected.
    \item Attribution: verified directly against the live
    API with authenticated and unauthenticated requests
    per role.
\end{itemize}

\subsection{Secondary evaluation metrics}
\begin{itemize}[leftmargin=0.7cm]
    \item \textbf{Explanation coverage:} percentage of
    delivered scores that have a non-empty, correctly
    normalised SHAP explanation (points sum to
    \texttt{score - SCORE\_MIN}).
    \item \textbf{Cold-start behaviour:} score assigned to
    a borrower with zero repayment events is exactly the
    configured prior (SHG bootstrap or independent base),
    with no model call.
    \item \textbf{Marketplace liveness:} lender
    accept/reject decisions on a loan request are reflected
    in both the lender's and the borrower's view without a
    stale read.
    \item \textbf{Lender reputation soundness:} a lender's
    computed reputation score responds monotonically to a
    lower offered interest rate relative to the platform
    average, given the same activity level.
\end{itemize}

\subsection{Criteria added after the payment-gateway and
lender-targeting iteration %
\BvrHint{\ldots problem / solution / evaluation criterion}}

\textbf{Problem:} Borrowers had no way to build real credit
history through the app -- repayment history only ever came
from offline synthetic data generation, never from a live
borrower action.\\
\textbf{Solution:} A live repayment flow (wallet balance +
dummy payment gateway) that pays down a real EMI.\\
\textbf{Evaluation criterion:} A successful payment changes
the borrower's credit score within the same request that
processes it -- verified on a real account: score moved
from 596 (Medium risk) to 886 (Very Low risk) once enough
on-time repayment history accumulated. A wallet balance
insufficient to cover the borrower's \emph{total} dues this
cycle blocks payment of every active loan, not just the one
clicked.

\medskip
\textbf{Problem:} A loan request always broadcast to every
eligible lender, with no way to direct it at one preferred
lender.\\
\textbf{Solution:} Optional lender targeting on a request,
with a retarget path if declined.\\
\textbf{Evaluation criterion:} A targeted request is visible
only in the targeted lender's queue (verified: other
eligible lenders do not see it), and a declined targeted
request can be retargeted rather than becoming a dead end.

\medskip
\textbf{Problem:} \texttt{GET /api/shgs} (the SHG browse list
used by both the lender and admin dashboards) computed each
SHG's aggregate repayment/attendance rate as a Python
property that walked \texttt{members -> loans ->
repayment\_events} per SHG -- an N+1 access pattern that, for
51 SHGs, materialised on the order of 25{,}000+
\texttt{RepaymentEvent} ORM objects on every request.\\
\textbf{Solution:} Replaced the per-SHG Python walk with
three \texttt{GROUP BY} SQL aggregate queries computed once
for all SHGs, joined back to each row by id.\\
\textbf{Evaluation criterion (measured, not estimated):}
median response time for \texttt{GET /api/shgs} against the
live dataset (51 SHGs, 1{,}079 individuals) dropped from
1{,}101--1{,}322\,ms to 14.5--26.7\,ms across repeated runs
-- roughly a 75$\times$ improvement -- with byte-for-byte
identical response content verified before and after.

\subsection{Pilot validation plan %
\BvrHint{\ldots high-level}}
\begin{itemize}[leftmargin=0.7cm]
    \item Exercise the full pipeline against realistic
    synthetic data (over a thousand individuals, dozens of
    SHGs, multiple lenders across several districts).
    \item Validate each role's dashboard end-to-end in a
    live browser session (not just unit tests) --
    screenshotting and checking console/network errors
    for each of the four logins.
\end{itemize}

\section{Scalability %
\BvrHint{\ldots with foresight}}
The proposition should scale in both theoretical and
practical senses.

\begin{enumerate}[leftmargin=0.7cm]
    \item \textbf{Scalable (theoretically):} the system
    supports growth in individuals, SHGs, and lenders by:
    \begin{itemize}[leftmargin=0.7cm]
        \item decoupling frontend and backend behind a
        REST API,
        \item indexing every foreign key (a gap we found
        and fixed during the build),
        \item paginating and server-side filtering list
        endpoints instead of client-side filtering over one
        page.
    \end{itemize}

    \item \textbf{Operationally deployable (practically):}
    the system runs on common web hosting:
    \begin{itemize}[leftmargin=0.7cm]
        \item a one-line environment-variable swap from
        SQLite to a managed Postgres instance,
        \item stateless bearer-token sessions rather than
        server-side sticky state,
        \item CI/CD pipeline for repeatable
        build/deploy.
    \end{itemize}
\end{enumerate}

\section{Engine availability heuristic}
A mature set of ``engines'' (tooling/services) is
available to implement this project:
\begin{itemize}[leftmargin=0.7cm]
    \item FastAPI for REST APIs, with automatic OpenAPI
    docs.
    \item SQLAlchemy ORM over SQLite/Postgres.
    \item XGBoost and SHAP for the scoring model and its
    explanations.
    \item scikit-learn (k-means, isolation forest) for
    geo-clustering and anomaly detection.
    \item React + Vite for the frontend.
    \item A token-based session-auth pattern (no new
    dependency needed for a demo-scale system).
\end{itemize}
We use these mature components to focus on pipeline
correctness, explainability, and access control rather
than inventing new infrastructure.

\section{Project scope and deliverables %
\BvrHint{\ldots iteration-friendly}}
\subsection{Initial deliverable (shipped) %
\BvrHint{\ldots first iterations}}
\begin{itemize}[leftmargin=0.7cm]
    \item Synthetic data generation (individuals, SHGs,
    lenders, loans, repayments, savings).
    \item Trained XGBoost + SHAP scoring engine, with
    cold-start bootstrap blending for both borrower paths.
    \item SHG-lender linkage, lender matching, and a loan
    offer/request marketplace.
    \item k-means district clustering and isolation-forest
    anomaly detection.
    \item A first working frontend across all of the above.
\end{itemize}

\subsection{Current iteration (prototype stage)}
\begin{itemize}[leftmargin=0.7cm]
    \item Four separate, role-gated logins (Borrower,
    Lender, SHG, Admin) replacing the previously-open
    single interface.
    \item Server-side authorization enforced on every
    endpoint, scoped to the authenticated caller's own
    identity.
    \item A lender reputation/incentive score rewarding
    competitive interest rates, informed by real-world
    analogues (RBI priority-sector-lending and co-lending
    incentive structures).
    \item A redesigned landing page surfacing real,
    data-derived borrower-improvement stories, SHG/lender
    spotlights, and district-level score trends.
\end{itemize}

\subsection{Subsequent deliverables}
\begin{itemize}[leftmargin=0.7cm]
    \item Periodic re-scoring so genuine score-history
    trends accumulate over time (currently each individual
    has a small number of scoring snapshots, since the
    pipeline has only been run a handful of times).
    \item Admin account-management/moderation actions
    beyond the current read-only oversight view.
    \item Real password hashing/short-lived signed tokens
    in place of the deliberately-minimal demo bearer-token
    scheme.
\end{itemize}

\section{Risks and mitigations}
\begin{itemize}[leftmargin=0.7cm]
    \item \textbf{Cold-start data scarcity:} mitigated by
    an explicit bootstrap-to-model blend rather than
    forcing the model to score borrowers it has no signal
    on.
    \item \textbf{Fairness of the independent-vs-SHG base
    score gap:} this is a stated policy parameter
    (\texttt{INDEPENDENT\_BASE\_SCORE} vs. the SHG
    bootstrap formula), not something the data derived --
    flagged explicitly rather than hidden, and a candidate
    for a sensitivity analysis.
    \item \textbf{Demo-data limitations for time-series
    features:} the synthetic dataset is scored in discrete
    batch runs, so "score improved over time" features are
    only as rich as the number of real scoring passes that
    have actually been run -- we surface real deltas only,
    never fabricated ones, even when that means a feature
    section is temporarily sparse.
    \item \textbf{Security of role-based auth:} mitigated
    by enforcing authorization server-side (401/403) on
    every endpoint, not relying on the frontend to hide
    UI elements.
\end{itemize}

\section{Summary}
This project proposes a deployable, explainable, dual-path
credit-scoring and lending-marketplace platform for
microfinance in India. It meets the course criteria by
delivering value iteratively and verifying each stage
end-to-end against real data, using CI/CD to support
frequent safe releases, and defining measurable evaluation
criteria (access-control correctness, explanation
coverage, cold-start correctness, and marketplace
liveness). It remains focused on engineering pipeline
correctness, explainability, and access control rather
than algorithmic novelty.

\end{document}
